% !TEX TS-program = xelatex
% !TEX encoding = UTF-8 Unicode
% !Mode:: "TeX:UTF-8"

\documentclass{resume}
\usepackage{zh_CN-Adobefonts_external} % Simplified Chinese Support using external fonts (./fonts/zh_CN-Adobe/)
%\usepackage{zh_CN-Adobefonts_internal} % Simplified Chinese Support using system fonts
\usepackage{linespacing_fix} % disable extra space before next section
\usepackage{cite}

\begin{document}
\pagenumbering{gobble} % suppress displaying page number

\name{李凯}

% {E-mail}{mobilephone}{homepage}
% be careful of _ in emaill address
\contactInfo{15301959284}{likai2023@126.com}{嵌入式工程师}{}
% {E-mail}{mobilephone}
% keep the last empty braces!
%\contactInfo{xxx@yuanbin.me}{(+86) 131-221-87xxx}{} 
\section{个人总结}
本人在校期间成绩优良，动手能力突出，善于系统思考与工程实践。对嵌入式系统开发、底层通信协议、上位机软件设计具有浓厚兴趣，注重技术细节与实际效果，具备较强的问题分析与解决能力。工作态度认真负责，自我驱动力强，乐于挑战新技术，追求高质量的工程落地与持续优化。

% \section{\faGraduationCap\ 教育背景}
\section{教育背景}
% 主修学位（按时间倒序）  
\datedsubsection{\textbf{上海大学}, 控制工程, \textit{在读硕士研究生}}{2023.9 -- 至今}  
\ \textbf{研究方向}: 嵌入式系统开发、工业控制算法  
\ \textbf{核心课程}: 嵌入式系统设计、自动控制原理、现代传感器技术、机器人控制  
% 若有奖学金、竞赛奖项等可补充  

\datedsubsection{\textbf{上海理工大学}, 机械设计及其自动化, \textit{工学学士}}{2019.9 -- 2023.6}  
\ \textbf{排名}: 前15\% 
\ \textbf{相关课程}: C语言程序设计、单片机原理与应用、数字电路设计、机电一体化系统  
% 补充与嵌入式相关的课程或项目  

% \section{\faCogs\ IT 技能}
\section{技术能力}
% increase linespacing [parsep=0.5ex]
\datedsubsection{\textbf{编程语言}}{}  
\ \ C/C++、Python、MATLAB、Java

\datedsubsection{\textbf{嵌入式开发}}{}  
\ \ 单片机开发：STM32、NXP S32K（汽车电子方向）、ESP32 
\ \ 嵌入式操作系统：Linux驱动基础与应用开发、Shell脚本
\ \ 硬件工具：Altium Designer、示波器/逻辑分析仪调试  

\datedsubsection{\textbf{其他工具}}{}  
\ \ Git版本控制、Keil/IAR集成开发环境、Qt、SolidWorks（机械设计基础）  

% \end{itemize}

\section{项目经历}
% 项目1：物联网设备开发（强相关）  
\datedsubsection{\textbf{自助式加注机支付系统（MCI2.0）} | 嵌入式物联网设备}{2023.10 -- 2024.6}  
\ \textbf{技术栈}: STM32 + TCP/IP协议栈 + IIC,SPI,UART,1-Wire + 传感器数据融合  
\ \begin{itemize}  
	\item 基于实时中断架构设计：  
	(1) 配置RTC实时时钟中断、外部GPIO中断（按键/传感器触发）、定时器中断  
	(2) 主循环事件驱动（状态机模式），中断服务例程（ISR）处理紧急任务（如电池低电压报警）   
	(3) 通过4G模块与服务器通信（MQTT协议（休眠模式+中断唤醒）  
	(4) 支持OTA远程升级，减少现场维护成本 
	
\end{itemize}  

% 项目2：嵌入式系统综合应用（突出软件能力，我更希望向着软件方向发展）  
% 项目2：嵌入式系统综合应用（突出软件能力，更偏向软件方向）  
\datedsubsection{\textbf{无人机智能巡检系统} | Android地面站软件与边缘计算协同}{2024.9 -- 至今}  
\begin{itemize}  
	\item \textbf{技术栈}: Android Studio (Kotlin/Java) + MAVLink v2.0 + FFmpeg + Mapbox + STM32F103/RK3588  
	\item \textbf{Android地面站核心功能开发}:  
	\begin{itemize}  
		\item 基于\textbf{MAVLink v2.0}协议实现飞控通信，解析HM30数传模块的实时数据（姿态、GPS、电池状态），封装控制指令；  
		\item 设计\textbf{多线程协同框架}：  
		\begin{itemize}  
			\item 主线程管理UI渲染及Mapbox路径显示；  
			\item 独立线程处理MAVLink协议解析/封装，保障飞控数据实时性；  
			\item 异步线程拉取RK3588的RTSP视频流，实现视频流与飞控数据的帧同步叠加显示；  
		\end{itemize}  
		\item 开发蓝牙透传模块，连接HM30数传硬件，支持平板直接控制无人机启停、巡航模式切换；  
		\item 集成边缘计算结果：接收RK3588的YOLOv5s缺陷检测结果，在Mapbox地图动态标注异常位置，并触发无人机悬停指令.  
	\end{itemize}  
\end{itemize}  

\section{实习经历}
\datedsubsection{\textbf{上海易沐科技有限公司}, 嵌入式工程师}{2024.08-2025.05}
\begin{itemize}
	\item \textbf{S32K148 Bootloader升级机制优化}
	\begin{itemize}
		\item 基于NXP S32K148单片机，整体重构Bootloader架构，设计并优化串口与CAN双通道IAP（在应用编程）升级机制。
		\item 改进升级协议，串口使用DMA接收，CAN采用FIFO中断机制，提升传输效率；将传统.srec文件升级切换为bin文件，整体升级时间由20分钟缩短至1分钟以内（固件大小700KB，波特率115200bps）。
		\item 完善升级跳转流程，确保升级过程稳定可靠，解决EVCC与SECC设备固件更新慢、易出错的问题，目前已应用于小批量试用项目中。
	\end{itemize}
	
	\item \textbf{Qt平台集成开发与优化}
	\begin{itemize}
		\item 主导开发Qt5.12.2集成平台，集成固件升级、工装检测与参数配置功能，提升生产测试效率。
		\item 使用Qt串口与CAN模块，支持固件一键升级及参数读写，设计UI界面实现人机交互操作，支持单独配置与批量导入导出功能。
		\item 搭建工装检测模块，支持人工触发检测流程，简化工装设备状态检测操作。
	\end{itemize}
	
	\item \textbf{交流桩环形充电算法设计}
	\begin{itemize}
		\item 基于图论理论，设计交流桩三环拓扑结构下的环形充电算法，优化充电模块功率分配，实现负载均衡。
		\item 支持节点间充电路径动态调整，提升系统可靠性，目前处于功能测试阶段。
	\end{itemize}
\end{itemize}



% \begin{onehalfspacing}
% \end{onehalfspacing}

% \datedsubsection{\textbf{DID-ACTE} 荷兰莱顿}{2015年}
% \role{本科毕业设计}{LIACS 交换生}
% 利用结巴分词对中国古文进行分词与词性标注，用已有领域知识训练形成 classifier 并对结果进行调优
% \begin{onehalfspacing}
% \begin{itemize}
%   \item 利用结巴分词对中国古文进行分词与词性标注
%   \item 利用已有领域知识训练形成 classifier, 并用分词结果进行测试反馈
%   \item 尝试不同规则，对 classifier 进行调优
% \end{itemize}
% \end{onehalfspacing}


% \section{\faHeartO\ 项目/作品摘要}
% \section{项目/作品摘要}
% \datedline{\textit{An Integrated Version of Security Monitor Vis System}, https://hijiangtao.github.io/ss-vis-component/ }{}
% \datedline{\textit{Dark-Tech}, https://github.com/hijiangtao/dark-tech/ }{}
% \datedline{\textit{融合社交网络数据挖掘的电视节目可视化分析系统}, https://hijiangtao.github.io/variety-show-hot-spot-vis/}{}
% \datedline{\textit{LeetCodeOJ Solutions}, https://github.com/hijiangtao/LeetCodeOJ}{}
% \datedline{\textit{Info-Vis}, https://github.com/ISCAS-VIS/infovis-ucas}{}


%% Reference
%\newpage
%\bibliographystyle{IEEETran}
%\bibliography{mycite}
\end{document}
